% ========================================================
% BÀI TẬP TOÁN 5 - KẾT NỐI TRI THỨC VỚI CUỘC SỐNG
% Bài 75: Ôn tập chung
% Định dạng: BTTeX Part III - Trắc nghiệm điền số ngắn
% ========================================================

% Câu 1
\begin{ex}
  Tính giá trị biểu thức: $2\,024 - 1\,000 : 8$. Kết quả là bao nhiêu?
  \shortans{1899}
  \loigiai{
    $2\,024 - 125 = 1\,899$.
  }
\end{ex}

% Câu 2
\begin{ex}
  Tính: $3{,}5 \times 12 + 6{,}5 \times 12$. Kết quả là bao nhiêu?
  \shortans{120}
  \loigiai{
    $(3{,}5 + 6{,}5) \times 12 = 10 \times 12 = 120$.
  }
\end{ex}

% Câu 3
\begin{ex}
  Tính: $\frac{3}{5} + \frac{2}{5} \times \frac{5}{4}$. Kết quả viết dưới dạng phân số tối giản có mẫu số là bao nhiêu?
  \shortans{10}
  \loigiai{
    $\frac{3}{5} + \frac{1}{2} = \frac{6}{10} + \frac{5}{10} = \frac{11}{10}$. Mẫu số là 10.
  }
\end{ex}

% Câu 4
\begin{ex}
  Ở biểu thức trên, tử số là bao nhiêu?
  \shortans{11}
  \loigiai{
    Tử số là 11.
  }
\end{ex}

% Câu 5
\begin{ex}
  Tìm $x$: $x \times 2{,}5 = 50$. Giá trị của $x$ là bao nhiêu?
  \shortans{20}
  \loigiai{
    $x = 50 : 2{,}5 = 20$.
  }
\end{ex}

% Câu 6
\begin{ex}
  Một hình chữ nhật có chiều dài $24\text{ m}$, chiều rộng bằng $\frac{2}{3}$ chiều dài. Chiều rộng là bao nhiêu mét?
  \shortans{16}
  \loigiai{
    $24 \times \frac{2}{3} = 16\text{ m}$.
  }
\end{ex}

% Câu 7
\begin{ex}
  Ở hình chữ nhật trên, diện tích là bao nhiêu mét vuông?
  \shortans{384}
  \loigiai{
    $24 \times 16 = 384\text{ m}^2$.
  }
\end{ex}

% Câu 8
\begin{ex}
  Một hình thang có hai đáy dài $18\text{ cm}$ và $12\text{ cm}$, chiều cao $10\text{ cm}$. Diện tích hình thang là bao nhiêu xăng-ti-mét vuông?
  \shortans{150}
  \loigiai{
    $(18 + 12) \times 10 : 2 = 150\text{ cm}^2$.
  }
\end{ex}

% Câu 9
\begin{ex}
  Một hình lập phương có diện tích toàn phần là $150\text{ cm}^2$. Thể tích hình lập phương là bao nhiêu xăng-ti-mét khối?
  \shortans{125}
  \loigiai{
    Diện tích 1 mặt: $150 : 6 = 25\text{ cm}^2$, cạnh $5\text{ cm}$. Thể tích: $5 \times 5 \times 5 = 125\text{ cm}^3$.
  }
\end{ex}

% Câu 10
\begin{ex}
  Một bể bơi dạng hình hộp chữ nhật dài $25\text{ m}$, rộng $10\text{ m}$, sâu $2\text{ m}$. Thể tích bể bơi là bao nhiêu mét khối?
  \shortans{500}
  \loigiai{
    $25 \times 10 \times 2 = 500\text{ m}^3$.
  }
\end{ex}

% Câu 11
\begin{ex}
  Quãng đường $AB$ dài $120\text{ km}$. Một ô tô đi từ $A$ đến $B$ với vận tốc $50\text{ km/h}$, đồng thời một xe máy đi từ $B$ về $A$ với vận tốc $30\text{ km/h}$. Sau bao nhiêu phút hai xe gặp nhau?
  \shortans{90}
  \loigiai{
    Tổng vận tốc: $80\text{ km/h}$. Thời gian gặp: $120 : 80 = 1{,}5\text{ giờ} = 90$ phút.
  }
\end{ex}

% Câu 12
\begin{ex}
  Một ca nô đi xuôi dòng với vận tốc $32\text{ km/h}$, ngược dòng $26\text{ km/h}$. Vận tốc của dòng nước là bao nhiêu km/h?
  \shortans{3}
  \loigiai{
    $(32 - 26) : 2 = 3\text{ km/h}$.
  }
\end{ex}

% Câu 13
\begin{ex}
  Ở bài trên, vận tốc thực của ca nô là bao nhiêu km/h?
  \shortans{29}
  \loigiai{
    $(32 + 26) : 2 = 29\text{ km/h}$.
  }
\end{ex}

% Câu 14
\begin{ex}
  Tìm hai số có tổng là 150 và tỉ số là $\frac{2}{3}$. Số lớn là bao nhiêu?
  \shortans{90}
  \loigiai{
    Tổng số phần: $2 + 3 = 5$. Số lớn: $(150 : 5) \times 3 = 90$.
  }
\end{ex}

% Câu 15
\begin{ex}
  Ở bài trên, số bé là bao nhiêu?
  \shortans{60}
  \loigiai{
    $(150 : 5) \times 2 = 60$.
  }
\end{ex}

% Câu 16
\begin{ex}
  Một lớp học có 40 học sinh, trong đó có 26 bạn nữ. Tỉ số phần trăm số bạn nữ là bao nhiêu phần trăm?
  \shortans{65}
  \loigiai{
    $(26 : 40) \times 100\% = 65\%$.
  }
\end{ex}

% Câu 17
\begin{ex}
  Giá niêm yết một đôi giày là $500\,000$ đồng, được giảm giá $20\%$. Số tiền được giảm là bao nhiêu nghìn đồng?
  \shortans{100}
  \loigiai{
    $500 \times 20\% = 100$ nghìn đồng.
  }
\end{ex}

% Câu 18
\begin{ex}
  Ở bài trên, khách hàng phải trả bao nhiêu nghìn đồng?
  \shortans{400}
  \loigiai{
    $500 - 100 = 400$ nghìn đồng.
  }
\end{ex}

% Câu 19
\begin{ex}
  Đổi $3\text{ tấn } 5\text{ tạ}$ ra ki-lô-gam được bao nhiêu ki-lô-gam?
  \shortans{3500}
  \loigiai{
    $3\,000 + 500 = 3\,500\text{ kg}$.
  }
\end{ex}

% Câu 20
\begin{ex}
  Đổi $2{,}5\text{ ha}$ ra mét vuông được bao nhiêu mét vuông?
  \shortans{25000}
  \loigiai{
    $2{,}5 \times 10\,000 = 25\,000\text{ m}^2$.
  }
\end{ex}

% Câu 21
\begin{ex}
  Đổi $3\text{ giờ } 45\text{ phút}$ ra phút được bao nhiêu phút?
  \shortans{225}
  \loigiai{
    $3 \times 60 + 45 = 225$ phút.
  }
\end{ex}

% Câu 22
\begin{ex}
  Tung xúc xắc 60 lần, mặt chẵn xuất hiện 36 lần. Tỉ số phần trăm mặt chẵn xuất hiện là bao nhiêu phần trăm?
  \shortans{60}
  \loigiai{
    $(36 : 60) \times 100\% = 60\%$.
  }
\end{ex}

% Câu 23
\begin{ex}
  Biểu đồ quạt biểu diễn 3 loại hoa trong vườn: Hoa hồng $40\%$, Hoa cúc $35\%$, Hoa lan $25\%$. Nếu có 400 bông hoa thì hoa lan có bao nhiêu bông?
  \shortans{100}
  \loigiai{
    $400 \times 25\% = 100$ bông.
  }
\end{ex}

% Câu 24
\begin{ex}
  Một người đi xe máy khởi hành lúc $7\text{ giờ } 30\text{ phút}$ và đến nơi lúc $9\text{ giờ } 15\text{ phút}$. Người đó đi hết bao nhiêu phút?
  \shortans{105}
  \loigiai{
    $9\text{h}15 - 7\text{h}30 = 1\text{ giờ } 45\text{ phút} = 105$ phút.
  }
\end{ex}

% Câu 25
\begin{ex}
  Một hình tròn có chu vi $62{,}8\text{ cm}$. Bán kính của hình tròn là bao nhiêu xăng-ti-mét (lấy $\pi = 3{,}14$)?
  \shortans{10}
  \loigiai{
    $62{,}8 : (2 \times 3{,}14) = 10\text{ cm}$.
  }
\end{ex}
